% Part: many-valued-logic
% Chapter: infinite-valued-logics
% Section: goedel

\documentclass[../../../include/open-logic-section]{subfiles}

\begin{document}

\olfileid[fr]{mvl}{inf}{god}

\olsection{Les logiques de G\"odel}

Les mêmes opérations définissent la matrice de G\"odel sur
les valeurs rationnelles et sur chacune des chaînes finies
$V_m$. La comparaison de ces matrices porte sur les
tautologies et sur la conséquence à partir de prémisses.

\begin{defn}\ollabel{def:goedel}
La logique de G\"odel à une infinité
de valeurs~$\LogGod[\infty]$ est
définie par la matrice suivante :
\begin{enumerate}
  \item Le langage propositionnel usuel
  $\Lang L_0$ avec $\lfalse$,
  $\lnot$, $\land$, $\lor$ et $\lif$.
  \item L'ensemble des valeurs de vérité
  $V_\infty$.
  \item $1$ est la seule valeur désignée,
  c'est-à-dire $V^+=\{1\}$.
  \item Les fonctions de vérité sont :
  \begin{align*}
    \tf{\lfalse} &=0\\
    \tf{\lnot}[\LogGod](x)
      &=\begin{cases}
        1 & \text{si }x=0,\\
        0 & \text{sinon,}
      \end{cases}\\
    \tf{\land}[\LogGod](x,y) &=\min(x,y)\\
    \tf{\lor}[\LogGod](x,y) &=\max(x,y)\\
    \tf{\lif}[\LogGod](x,y)
      &=\begin{cases}
        1 & \text{si }x\le y,\\
        y & \text{sinon.}
      \end{cases}
  \end{align*}
\end{enumerate}
La logique de G\"odel à $m$ valeurs
est définie de la même façon, à ceci
près que $V=V_m$.
\end{defn}

\begin{prop}
  La logique $\LogGod[3]$ définie
  par la \olref[thr][god]{defn:goedel}
  coïncide, sous l'identification
  $\True=1$, $\Undef=1/2$ et
  $\False=0$, avec la logique
  $\LogGod[3]$ définie par la
  \olref{def:goedel}.
\end{prop}

\begin{proof}
  Il suffit de comparer les tables
  de vérité des connecteurs données
  dans la \olref[thr][god]{defn:goedel}
  avec celles que déterminent les
  équations de la
  \olref{def:goedel} :
  \begin{center}
    \begin{tabular}{c|c}
      $\tf{\lnot}[\LogGod[3]]$ & \\
      \hline
      $1$ & $0$ \\
      $1/2$ & $0$ \\
      $0$ & $1$
    \end{tabular}
    \quad
    \begin{tabular}{c|ccc}
      $\tf{\land}[\LogGod]$ & $1$ & $1/2$ & $0$ \\
      \hline
      $1$ & $1$ & $1/2$ & $0$ \\
      $1/2$ & $1/2$ & $1/2$ & $0$ \\
      $0$ & $0$ & $0$ & $0$
    \end{tabular}
    \\[2ex]
    \begin{tabular}{c|ccc}
      $\tf{\lor}[\LogGod]$ & $1$ & $1/2$ & $0$ \\
      \hline
      $1$ & $1$ & $1$ & $1$ \\
      $1/2$ & $1$ & $1/2$ & $1/2$ \\
      $0$ & $1$ & $1/2$ & $0$
    \end{tabular}
    \quad
    \begin{tabular}{c|ccc}
      $\tf{\lif}[\LogGod]$ & $1$ & $1/2$ & $0$ \\
      \hline
      $1$ & $1$ & $1/2$ & $0$ \\
      $1/2$ & $1$ & $1$ & $0$ \\
      $0$ & $1$ & $1$ & $1$
    \end{tabular}
  \end{center}
\end{proof}

\begin{prop}\ollabel{prop:god-infty-m}
  Si $\Gamma\Entails[\LogGod[\infty]] !B$,
  alors
  $\Gamma\Entails[\LogGod[m]] !B$
  pour tout~$m\ge2$.
\end{prop}

\begin{proof}
  Exercice.
\end{proof}

\begin{prob}
  Démontrer la
  \olref[mvl][inf][god]{prop:god-infty-m}.
\end{prob}

La réciproque vaut si $\Gamma$ est fini. Supposons qu'une
valuation dans $V_\infty$ désigne toutes les prémisses et donne
à $!B$ une valeur inférieure à $1$. Les valeurs des variables
de $\Gamma$ et de $!B$ ont un dénominateur commun $d$.
Elles appartiennent à $V_{d+1}$. Comme les opérations de G\"odel
se restreignent à cette chaîne, on obtient le même contre-modèle
en donnant la valeur $0$ aux variables restantes.

Pour voir le rôle de la finitude, prenons des variables
$p_0,p_1,\ldots,q$ deux à deux distinctes et, pour tout $i\geq 0$,
la prémisse $(p_{i+1}\lif p_i)\lif q$. Si $q$ a une valeur
inférieure à $1$, chaque prémisse impose que la valeur de
$p_{i+1}$ dépasse strictement celle de $p_i$ et que celle de
$p_i$ soit au plus la valeur de $q$. Une chaîne finie ne
peut contenir une telle suite. Dans $V_\infty$, on peut donner
à $q$ la valeur $1/2$ et à $p_i$ la valeur $1/2-1/(i+2)$.
Toutes les prémisses valent alors $1$ et $q$ vaut $1/2$.

Comme $\LogGod[3]$,
$\LogGod[\infty]$ admet comme
tautologies toutes les formules
valides intuitionnistement. Les mêmes
exemples de non-tautologies pour
$\LogGod[3]$ restent des
non-tautologies pour
$\LogGod[\infty]$ :
\begin{align*}
  & p\lor\lnot p && (p\lif q)\lif(\lnot p\lor q) \\
  & \lnot\lnot p\lif p && \lnot(\lnot p\land\lnot q)\lif(p\lor q) \\
  & ((p\lif q)\lif p)\lif p && \lnot(p\lif q)\lif(p\land\lnot q)
\end{align*}
L'exemple d'une formule non valide
intuitionnistement mais tautologique
dans~$\LogGod[3]$,
$(p\lif q)\lor(q\lif p)$, est aussi
une tautologie de~$\LogGod[\infty]$.
En fait, l'ensemble des tautologies
de~$\LogGod[\infty]$ est celui de la
logique intuitionniste augmentée du
schéma
$(!A\lif !B)\lor(!B\lif !A)$.
Michael Dummett l'a montré; c'est
pourquoi $\LogGod[\infty]$ est
souvent appelée \emph{logique de
G\"odel--Dummett}~$\Log{LC}$.

\begin{prob}
  Montrer que
  $(p\lif q)\lor(q\lif p)$ est une
  tautologie de~$\LogGod[\infty]$.
\end{prob}

\begin{prob}
  Montrer que
  $(p\lif q)\lor(q\lif r)\lor(r\lif s)$,
  qui est une tautologie de
  $\LogGod[3]$, n'est pas une
  tautologie de~$\LogGod[\infty]$.
\end{prob}

\end{document}
